This work studies a failure of label-free training that makes models answer questions they cannot answer. Answering a defective question with one committed number can harm users wherever the defect goes unnoticed: a tutoring system that answers an ill-posed exercise tells a student that the exercise has a unique answer; a scientific or analytical tool that silently fills in a missing assumption returns a number whose basis the user never sees; and an assistant in a customer-facing setting may settle an incomplete or ambiguous request by guessing, without saying so. The risk is smaller when the answer states its assumption or declines, as half to two thirds of our trained solvers' samples do, and largest where only the final number is read. A test-time abstention option, which works in our experiments, may be insufficient on its own, because users may ignore the abstentions of a model that usually answers. We therefore recommend both measuring well-posedness independently of the model's own agreement during training and offering an explicit abstention option at test time. The released curricula can also be misused: their problems are written by models, and their labels come from language-model judges, with human labels for only 240 problems, so training or evaluating on them as if the labels were ground truth would propagate the judges' errors; they should not be used as ground truth for answerability without human review. On well-posed problems, the judges' false-positive rate is about 5\%$^{\dagger}$ on intact planted problems and higher in later rounds on the human-labelled problems (blind Opus: 2.9\% at round one, 21--23.5\% at rounds three and five; Figure~\ref{fig:human}), and the release marks the source of every label.

The training data consist of problems written by open models and of premise-deleted and intact GSM-Symbolic variants; ill-posedness labels come from language-model judges (Claude Opus, with Qwen2.5-32B-Instruct as a check) and, for 240 problems, from three external mathematicians (two PhD students and a postdoctoral researcher), who were paid USD 500 each for about five hours of work, consented to the use and anonymised release of their labels, and saw no personal data. The sample of Appendix~\ref{app:readings} was labelled by two external annotators on the same terms of consent and data, for USD 1,000 each and 20.5 and 18.5 hours of work. The annotation was not reviewed by an ethics board. We use UMWP (CC BY-SA 4.0), SUM, MiP, GSM8K and MATH (MIT) and GSM-Symbolic, whose variants we release only as patches to respect its licence, and Qwen3, Qwen2.5, Phi-4-mini (MIT) and Llama 3.2 models under their licences.
